\section{Fundamentos de Aprendizaje Automático y Aprendizaje Profundo}

El desarrollo de sistemas modernos de gestión vehicular en estacionamientos ha
migrado desde métodos deterministas basados en procesamiento de imágenes clásico
hacia arquitecturas de aprendizaje profundo. Esta transición permite abordar la
alta variabilidad de los entornos de captura reales (como cambios de iluminación,
ángulos de inclinación de la cámara y oclusiones parciales) mediante modelos
capaces de extraer jerarquías de características discriminativas a partir de
grandes volúmenes de datos.

\subsection{El Paradigma del Aprendizaje Supervisado y Partición de Datos}

Como señalan Russell y Norvig \cite{russell2020artificial}, el aprendizaje
supervisado constituye un paradigma donde un algoritmo deduce una función de
mapeo a partir de pares de entrenamiento formados por muestras de entrada y sus
correspondientes etiquetas de referencia. En el ámbito del control de acceso
vehicular, cada entrada corresponde a un fotograma extraído de las cámaras
instaladas en las barreras del parqueadero, mientras que las etiquetas
representan el tipo de vehículo, la categoría de color o las coordenadas y
contenido alfanumérico de la matrícula.

Para garantizar que el modelo adquiera capacidad predictiva sobre instancias no
observadas previamente y no memorice las placas de los automóviles de
entrenamiento, Mohri et al. \cite{mohri2018foundations} y Goodfellow et al.
\cite{goodfellow2016deep} establecen la necesidad de aplicar una partición de
datos en tres conjuntos independientes:

\begin{itemize}
    \item \textbf{Conjunto de Entrenamiento (\textit{Training Set}):} Utilizado
        por el algoritmo para optimizar iterativamente los pesos internos del
        modelo a partir de miles de capturas de vehículos en el punto de acceso.

    \item \textbf{Conjunto de Validación (\textit{Validation Set}):} Destinado a
        monitorear el desempeño durante las etapas de entrenamiento, ajustar
        hiperparámetros (como la tasa de aprendizaje o el tamaño de filtro) y
        seleccionar la mejor arquitectura sin sesgar la evaluación.

    \item \textbf{Conjunto de Prueba (\textit{Test Set}):} Reservado
        estrictamente para medir la precisión final de generalización ante
        vehículos y matrículas nunca antes procesados por el sistema, simulando
        la operación real del estacionamiento.
\end{itemize}

\subsection{Arquitectura Base: Neuronas, Propagación y Activaciones no Lineales}

De acuerdo con Bishop y Bishop \cite{bishop2024deep}, las redes neuronales
alimentadas hacia adelante (\textit{feedforward networks}) procesan la
información de manera unidireccional a través de capas compuestas por unidades
computacionales (neuronas). En el nivel básico, cada neurona efectúa una
combinación lineal de las señales de entrada mediante pesos y sesgos, pasando el
resultado a través de una función de activación.

El propósito fundamental de incorporar funciones de activación no lineales (como
ReLU o SiLU) es romper la linealidad del sistema. Sin estas transformaciones,
una red de múltiples capas colapsaría matemáticamente en una simple
transformación afín, siendo incapaz de resolver fronteras de decisión complejas.
En la percepción vehicular, la no linealidad es la que permite al sistema
proyectar las intensidades de los píxeles hacia espacios abstractos donde sea
posible, por ejemplo, separar el color cromático de la carrocería de un
automóvil respecto a las sombras proyectadas en el pavimento o las reflexiones
del entorno.

\subsection{Optimización de Modelos: Descenso de Gradiente y Funciones de Pérdida}

El proceso de entrenamiento se formula como la minimización de una función de
pérdida, la cual cuantifica la discrepancia entre las predicciones emitidas por
el modelo y las anotaciones reales. Para ajustar los parámetros de la red, los
algoritmos de optimización —como el Descenso de Gradiente Estocástico (SGD) o
estimadores adaptativos como Adam \cite{goodfellow2016deep, bishop2024deep}—
calculan el gradiente de la función de pérdida con respecto a cada peso mediante
el algoritmo de retropropagación (\textit{backpropagation}), que aplica la regla
de la cadena en sentido inverso a través de la arquitectura.

En el control de acceso vehicular, el proceso de optimización persigue un
equilibrio operativo:

\begin{itemize}
    \item \textbf{Regresión de cajas delimitadoras:} Minimiza el error de
        localización geométrica entre la caja predicha por el detector y la
        posición real de la placa o del vehículo.

    \item \textbf{Minimización de falsos positivos:} Reduce la probabilidad de
        clasificar elementos del fondo (como rejas, sombras o señalética de
        tráfico) como vehículos autorizados o matrículas válidas.

\end{itemize}

\subsection{Generalización: Sesgo-Varianza, Sobreajuste y Regularización}

El desempeño de un modelo en producción está gobernado por el dilema entre el
sesgo (incapacidad del modelo para capturar la estructura subyacente de los
datos) y la varianza (sensibilidad excesiva del modelo a las fluctuaciones del
conjunto de entrenamiento). Un modelo con alto sesgo incurre en subajuste
(\textit{underfitting}), mientras que una varianza alta deriva en sobreajuste
(\textit{overfitting}).

En el contexto de la gestión de parqueaderos, el sobreajuste representa la
principal causa de falla: ocurre cuando la red neuronal ``memoriza'' las
condiciones particulares de la sesión de entrenamiento —por ejemplo, la
iluminación solar directa de la mañana o la orientación fija de un grupo de
vehículos locales—. Al desplegar el sistema en un entorno real dinámico, el
modelo sobreajustado colapsa al enfrentar cambios de iluminación pasante
(nubosidad o noche), lluvia sobre la calzada o variaciones en la trayectoria del
automóvil al aproximarse a la barrera.

Para mitigar este problema y asegurar una alta capacidad de generalización,
Srivastava et al. \cite{srivastava2014dropout} y LeCun et al. 
\cite{lecun2015deep} proponen técnicas de regularización estructurales y
estocásticas:

\begin{enumerate}
    \item \textbf{Dropout:} Desactiva aleatoriamente una fracción de neuronas
        durante cada iteración de entrenamiento, impidiendo la co-adaptación de
        características y obligando a la red a aprender representaciones
        redundantes y robustas \cite{srivastava2014dropout}.
        
    \item \textbf{Penalizaciones de Norma ($L_2$ / \textit{Weight Decay}):}
        Restringen la magnitud de los pesos ponderados, suavizando la función de
        decisión para que pequeñas variaciones en los píxeles de entrada no
        alteren drásticamente la predicción.

    \item \textbf{Normalización por Lotes (\textit{Batch Normalization}):}
        Estabiliza la distribución de las activaciones intermedias, acelerando la
        convergencia y actuando como un regularizador secundario.
\end{enumerate}

\subsection{Redes Neuronales Convolucionales (CNN) en la Percepción Visual}

Como establecen LeCun et al. \cite{lecun2015deep} y Szeliski
\cite{szeliski2022computer}, las Redes Neuronales Convolucionales (CNN) están
optimizadas para datos con topología en rejilla bidimensional mediante tres
mecanismos: campos receptivos locales, compartición de pesos y submuestreo
(\textit{pooling}). Estas propiedades otorgan invariancia a la traslación y
reducen radicalmente la cantidad de parámetros frente a capas densas
convencionales.

El procesamiento de una escena en el acceso al parqueadero se realiza mediante
la extracción jerárquica de características visuales:
\begin{equation}
    (I * K)(i, j) = \sum_{m} \sum_{n} I(i + m, j + n) K(m, n) \quad
\end{equation}
Donde $I$ es la imagen capturada por la cámara y $K$ es el filtro espacial.

\begin{itemize}
    \item \textbf{Capas Convolucionales Iniciales (Bajo Nivel):} Filtros
        espaciales pequeños reaccionan ante gradientes de intensidad y primitivas
        geométricas elementales, como los bordes metálicos de la carrocería, el
        contorno del parachoques o las líneas de alto contraste del marco de la
        placa de matrícula.

    \item \textbf{Capas Convolucionales Intermedias (Nivel Medio):} Los mapas de
        características agrupan las primitivas simples en estructuras compuestas,
        detectando la simetría rectangular de las placas, la forma circular de
        las ruedas o la textura de los faros.

    \item \textbf{Capas Convolucionales Profundas (Alto Nivel):} Sintetizan la
        información espacial en representaciones semánticas globales,
        permitiendo al sistema clasificar la silueta completa del vehículo 
        (sedán, camioneta, minibús o motocicleta) y distinguir la textura
        alfanumérica de los caracteres en la matrícula.

    \item \textbf{Capas de Submuestreo (\textit{Pooling}):} Reducen la resolución
        espacial de los mapas de características, reteniendo las respuestas más
        dominantes (\textit{Max-Pooling}) y otorgando tolerancia al sistema ante
        pequeñas variaciones en la posición del vehículo dentro del carril de
        ingreso.
\end{itemize}

\subsection{Aprendizaje por Transferencia (Transfer Learning), Fine-Tuning y
Aumento de Datos}

Entrenar un detector de objetos desde cero requiere millones de imágenes anotadas
y una enorme capacidad computacional. Para superar esta limitación en
aplicaciones del entorno real, Pan y Yang \cite{pan2010survey} formalizan el
Aprendizaje por Transferencia (\textit{Transfer Learning}), técnica que permite
transferir el conocimiento extraído de un dominio fuente masivo $\mathcal{D}_S$
(como los conjuntos de datos COCO u Objects365) hacia un dominio objetivo
$\mathcal{D}_T$.

Sustentado en las demostraciones empíricas de Yosinski et al.
\cite{yosinski2014transferable}, las capas iniciales de una CNN entrenada en
datasets generales capturan descriptores de bajo nivel universales. Por lo tanto,
en lugar de reiniciar los pesos, se adopta una arquitectura preentrenada de
vanguardia (como la familia YOLO) y se ejecuta un ajuste fino
(\textit{Fine-Tuning}), adaptando las capas superiores a las condiciones locales
del proyecto:

\begin{enumerate}
    \item \textbf{Especialización a la Morfología de Placas Locales:} Los modelos
        preentrenados globalmente están optimizados para formatos de matrículas
        estadounidenses o europeas. El \textit{fine-tuning} reconfigura los
        cabezales de predicción para detectar la geometría, color de fondo,
        relación de aspecto y tipografía estándar de las matrículas del contexto
        nacional (Bolivia).

    \item \textbf{Categorización de la Flota Vehicular Específica:} Permite
        adaptar la red para clasificar categorías de transporte particulares del
        entorno urbano local (como minibuses o ``surubíes'') que no existen de
        forma explícita en datasets internacionales.

    \item \textbf{Calibración a la Geometría de Captura:} Adapta el modelo a la
        distancia focal, ángulo de inclinación y altura fija de las cámaras
        instaladas en la infraestructura del parqueadero.
\end{enumerate}

De forma complementaria, el Aumento de Datos (\textit{Data Augmentation}) actúa
como un regularizador implícito que incrementa la variabilidad del conjunto de
entrenamiento sin costo de etiquetado adicional \cite{goodfellow2016deep,
bishop2024deep}. En el sistema de control de acceso, el aumento de datos simula
condiciones físicas y ambientales adversas:

\begin{itemize}
    \item \textbf{Variaciones Fotométricas (HSV y Brillo):} Simulan cambios
        drásticos de iluminación solar, reflejos o deslumbramientos sobre el
        parabrisas y condiciones de baja visibilidad nocturna.

    \item \textbf{Transformaciones Geométricas (Escalado, Rotación y
        Perspectiva):} Modulan las distintas trayectorias de ingreso de los
        automóviles, simulando maniobras descentradas o ángulos de aproximación
        no frontales.

    \item \textbf{Simulación de Intemperie y Oclusiones:} Inyecta artefactos
        sintéticos como ruido estocástico, desenfoque de movimiento y oclusiones
        parciales (\textit{Mosaic} / \textit{Mixup}) para representar gotas de
        lluvia sobre la lente, suciedad en la calzada o barro sobre la placa
        vehicular.
\end{itemize}

Mediante la integración del \textit{Transfer Learning}, el \textit{Fine-Tuning} y
el \textit{Data Augmentation}, el sistema alcanza la estabilidad operativa
requerida para garantizar un control de acceso automatizado, preciso y resistente
a las inclemencias del entorno real.
